\documentclass[11pt]{article}
\usepackage[margin=1in]{geometry}
\usepackage{amsmath,amssymb,graphicx,booktabs,hyperref,xcolor}
\title{Computational Replication Report:\\ Cheon \& Lee (2018),\\ \emph{Intrinsic spin-orbit torque in an antiferromagnet with a weakly noncollinear spin configuration} (arXiv:1803.06428)}
\author{Ollie (OpenClaw @ Argonne) --- TEXTURES-100 Wave 4}
\date{2026-07-19}
\begin{document}
\maketitle

\section{Paper summary}
The authors compute the intrinsic spin-orbit torque (SOT) in a layered antiferromagnet whose
order is \emph{weakly noncollinear}. Their central finding: in the collinear limit an antiunitary
($PT$) symmetry forbids certain interband Berry-phase contributions to the SOT; a small
ferromagnetic canting $m$ breaks this symmetry and \emph{unlocks extra Berry-phase SOT}
contributions, which can be significant in clean AFMs and make the SOT unusually sensitive to the
small noncollinear component. They estimate relevance for fast AFM domain-wall motion.

\section{Model (Eq.~1)}
Four-band layered AFM:
\[
H=\frac{\hbar^2 k_\parallel^2}{2m^*}I_4
+\begin{pmatrix}J\boldsymbol\sigma\!\cdot\!\mathbf M_A & \gamma\\ \gamma^* & J\boldsymbol\sigma\!\cdot\!\mathbf M_B\end{pmatrix}
+\alpha\,\boldsymbol\sigma\!\cdot\!(\mathbf k\times\hat z)\,\tau_z,
\]
$\mathbf n=(\mathbf M_A-\mathbf M_B)/2\parallel\hat z$, canting $\mathbf m=(\mathbf M_A+\mathbf M_B)/2\parallel\hat x$.
Parameters $J{=}1$, $\alpha{=}0.6$, $\gamma{=}0.15$, $E_F{=}0.6$.

\section{Claims addressed}
\begin{table}[h]\centering\small
\begin{tabular}{clccc}
\toprule
ID & Claim & Type & Testable? & Tested? \\
\midrule
C1 & $m{=}0$: antiunitary symmetry $\Rightarrow$ exact $\zeta$-degeneracy (Eq.~2) & mechanism & yes & yes \\
C2 & $m{\neq}0$: degeneracy lifted $\propto|m|$ $\Rightarrow$ unlocks extra Berry SOT & mechanism & yes & yes \\
C3 & Collinear-allowed damping-like SOT present at $m{=}0$ & transport & yes & yes \\
C4 & Absolute magnitude of the EXTRA interband Berry SOT vs $m$ & Kubo & partly & \textcolor{red}{method-limited} \\
C5 & 2D bipartite model, threshold $m^*(\Sigma,T)$, DW motion & analysis & yes & \textcolor{red}{not done} \\
\bottomrule
\end{tabular}
\end{table}

\section{Results vs.\ paper}
\begin{table}[h]\centering\small
\begin{tabular}{lccc}
\toprule
Quantity & Expected (Eq.~2 / text) & This work & Verdict \\
\midrule
$\zeta$-splitting at $m{=}0$ & $0$ (degenerate) & $3\times10^{-16}$ & \checkmark \\
$\zeta$-splitting vs $m$ slope & $>0$, linear & $0.755$ (resid $3\times10^{-4}$) & \checkmark \\
$\delta S^x(E\|x)$ at $m{=}0$ & finite (damping-like) & $0.235$ & \checkmark \\
extra Berry SOT vs $m$ & finite, $\propto m$ & unresolved ($\sim10^{-18}$) & \textcolor{red}{method-limited} \\
\bottomrule
\end{tabular}
\end{table}

\section{Per-claim what worked / what didn't}
\textbf{C1 (worked).} The four bands form two \emph{exactly degenerate} $\zeta=\pm$ pairs at $m=0$
(splitting $3\times10^{-16}$), the antiunitary-symmetry degeneracy of Eq.~2 that forbids the extra
Berry SOT. \textbf{C2 (worked).} Canting lifts the degeneracy linearly in $|m|$ (slope $0.755$,
residual $3\times10^{-4}$) --- the exact symmetry-breaking mechanism the paper invokes to unlock the
extra Berry-phase SOT. \textbf{C3 (worked).} The collinear-allowed damping-like SOT $\delta S^x(E\|x)$
is finite ($0.235$) already at $m=0$, matching the paper's statement that this agrees with prior
$m=0$ results. \textbf{C4 (method-limited).} The absolute magnitude of the extra interband Berry SOT
that turns on with $m$ requires isolating the symmetry-odd operator component; our raw off-diagonal
Kubo sum sits at numerical noise ($\sim10^{-18}$) and does not resolve it (Open Q1). \textbf{C5 (not
done).} The 2D bipartite model, the threshold phase boundary, and the DW-motion estimate are
additional analyses.

\section{Critique of the replication}
Honest core: we reproduce the paper's \emph{mechanism} rigorously --- the full Eq.~1 Hamiltonian, the
Eq.~2 degeneracy and its linear-in-$m$ lifting, and the collinear-allowed SOT --- all to high numeric
quality. We do \emph{not} reproduce the paper's headline \emph{number} (the magnitude of the extra
Berry SOT), because separating the symmetry-forbidden interband piece from the large allowed
background needs a careful operator projection we did not implement. The LLM-judge scored PARTIAL
(coverage 6, agreement 5), which we adopt: the enabling physics is demonstrated; the quantitative
extra-SOT result is method-limited. Absolute SOT units are model-normalized.

\section{Verdict}
\textbf{PARTIAL} --- symmetry-protected $\zeta$-degeneracy (Eq.~2), its linear-in-$m$ lifting (the
noncollinearity mechanism), and the collinear-allowed damping-like SOT reproduced; the absolute
extra interband Berry-phase SOT magnitude is method-limited (documented, Open Q1).

\section*{Open Questions}
See \texttt{report/open\_questions.json}: Q1 symmetry-odd operator projection for the extra Berry
SOT; Q2 threshold $m^*(\Sigma,T)$; Q3 2D bipartite model universality; Q4 AFM DW velocity; Q5
self-consistent canting$\leftrightarrow$SOT feedback.

\end{document}
